\subsection{赋值环的理想}

定义 2. 设 G 是有序集。若 G 的子集 M 满足：关系 \(x \in \mathbf{M}\) 与 \(y \geqslant x\) 蕴涵 \(y \in \mathbf{M}\)，则称 M 为优势集。

设 K 是域，v 是 K 上的赋值，A 是 v 的环，G 是 v 的阶群。对每个优势子集 \(M \subset G\)，用 \(\mathfrak{a}(M)\) 表示由 \(x \in K\) 中使 \(v(x) \in M \cup \{+\infty\}\) 的元素组成之集。显然 \(\mathfrak{a}(M)\) 是 K 的子 A-模。

命题 7. 映射 \(\mathbf{M} \mapsto \mathfrak{a}(\mathbf{M})\) 是从 \(\mathbf{G}\) 的优势子集所成之集到 \(\mathbf{K}\) 的子 A-模所成之集上的递增双射。

设 \(\mathfrak{b}\) 是 \(\mathbf{K}\) 的子 A-模。由诸 \(v(x)\)（\(x \in \mathfrak{b} - (0)\)）组成的集是优势子集 \(\mathbf{M}(\mathfrak{b})\)（\(\mathbf{G}\) 的）。若能证明下列等式，命题 7 即得证：

\begin{enumerate}
\def\labelenumi{(\arabic{enumi})}
\setcounter{enumi}{1}
\item
  \(\mathbf{M}(\mathfrak{a}(\mathbf{N})) = \mathbf{N}\)（对每个优势子集 \(\mathbf{N}\)，\(\mathbf{G}\) 中的）；
\item
  \(\mathfrak{a}(\mathbf{M}(\mathfrak{b})) = \mathfrak{b}\)（对每个子 A-模 \(\mathfrak{b}\)，\(\mathbf{K}\) 的）。
\end{enumerate}

公式 (2) 是容易的，因为对所有 \(m \in \mathbf{N}\)，存在 \(x \in \mathbf{K}\) 使 \(v(x) = m\)。于是显然 \(\mathfrak{b} \subset \mathfrak{a}(\mathbf{M}(\mathfrak{b}))\)；反之，设 \(x \in \mathfrak{a}(\mathbf{M}(\mathfrak{b}))\)，并设 \(x \neq 0\)；

则 \(v(x) \in \mathbf{M}(\mathfrak{b})\)，从而存在 \(y \in \mathfrak{b}\) 使 \(v(x) = v(y)\)；于是 \(x = uy\)，其中 \(v(u) = 0\)，这证明 \(x \in \mathrm{Ay} \subset \mathfrak{b}\)，证明完成。

推论. 设 \(G_{+}\) 是 G 中正元素所成之集。映射 \(\mathbf{M} \mapsto \mathfrak{a}(\mathbf{M})\) 是从 \(G_{+}\) 的优势子集所成之集到 A 的理想所成之集上的双射。

由于 \(\mathbf{A} = \mathfrak{a}(\mathbf{G}_{+}),\mathfrak{a}(\mathbf{M})\subset \mathbf{A}\) 等价于 \(\mathbf{M}\subset \mathbf{G}_{+}\)。

例如，极大理想 \(m(A)\) 等于 \(a(S)\)，其中 S 表示 G 的严格正元素所成之集。
% semantic-fragments: legacy-input-seam
\relax{}
